% Source: https://artofproblemsolving.com/community/c3h1996733_what_happens_now
% Stable identity: aops-topic-1996733-q01
\textbf{(a)} With $n=-3$,
\[
\binom{-3}{k}=\dfrac{(-3)(-4)\cdots(-3-k+1)}{k!}=(-1)^k\dfrac{3\cdot4\cdots(k+2)}{k!}=(-1)^k\dfrac{(k+2)!}{2\,k!}=(-1)^k\binom{k+2}{2}.
\]
Replacing $x$ by $-x$ cancels the sign: $(1-x)^{-3}=\sum_k\binom{k+2}{2}x^k$. This also has a counting meaning: $\dfrac{1}{(1-x)^3}=(1+x+x^2+\cdots)^3$, and the coefficient of $x^k$ is the number of ways to write $k=a+b+c$ with $a,b,c\ge0$, which is $\binom{k+2}{2}$ by stars and bars. For $k=4$ the coefficient is $\binom62=\boxed{15}$.

\textbf{(b)} With $n=\dfrac{1}{2}$:
\[
\binom{\dfrac{1}{2}}{1}=\dfrac{1}{2},\qquad \binom{\dfrac{1}{2}}{2}=\dfrac{\dfrac{1}{2}\cdot\left(-\dfrac{1}{2}\right)}{2}=-\dfrac{1}{8},\qquad \binom{\dfrac{1}{2}}{3}=\dfrac{\dfrac{1}{2}\cdot\left(-\dfrac{1}{2}\right)\cdot\left(-\dfrac{3}{2}\right)}{6}=\dfrac{1}{16}.
\]
So
\[
\sqrt{1+x}=1+\dfrac{x}{2}-\dfrac{x^2}{8}+\dfrac{x^3}{16}-\cdots .
\]
